\chapter{Final remarks and discussion}
\label{chap:final} In this thesis, we have presented a series of innovative
methods for rendering realistic human avatars using sparse conditioning
signals. We began by enhancing Neural Radiance Fields (NeRFs) with
controllability, introducing our \conerf model. Unlike concurrent approaches,
\conerf enables interpretable control over the 3D representation of a subject
while preserving the high quality of the output.

Recognizing that \conerf primarily focuses on spatially localized controls,
such as eye movements, we developed \blendfields. This method generates
animatable human head avatars from a single multi-view capture, predicting
realistic, expression-dependent skin deformations in a few-shot data regime.
By leveraging physically-based deformation features of underlying tetrahedra,
\blendfields accurately determines skin shading adjustments. Our experiments
demonstrated that \blendfields achieves impressive results from minimal data
and surpasses other methods reliant on learnable conditioning signals.

We then tackled the challenge of disentangling subject color from lighting
conditions in trained radiance fields. To address this, we introduced
\lumigauss, a model that learns lighting features purely from data. By
constraining the model so that the output color is a product of subject color
and lighting attributes, \lumigauss facilitates intuitive scene relighting
after training. Unlike methods such as~\textcite{saito2024relightable}, which
require extensive datasets of subject images taken under varied lighting
conditions, \lumigauss delivers high-fidelity results with minimal data.

Finally, we proposed a general approach to adapt 3D Gaussian Splats (3DGS)
representations to computational constraints. Our \clog method achieves this
by modeling 3D Gaussians on a 2D grid. We applied
PLAS~\cite{morgenstern2024compact} to reorganize the grid for compatibility
with frequency-based modulation. We then introduced a continuous modulator
that adjusts the grid based on conditioning signals, producing a constrained
representation. Our approach is among the first to enable continuous
representation adaptation during inference, without retraining, while
maintaining performance comparable to methods with fixed-scale modeling.

These contributions collectively address the research questions outlined
in~\cref{chap:introduction}. Together, they form a comprehensive framework for
radiance fields capable of: being controlled with a few annotation strokes,
animating human avatars with realistic, expression-dependent wrinkles,
relighting scenes using a small set of varied lighting images, adapting
representations to fit computational budgets. We see these advancements as
steps toward making NeRFs and 3DGS more accessible to the public and fostering
their proliferation in fields like 3D medical imaging.

\section{Future Work and Open Questions}
  Our research opens doors to several promising directions that warrant further
  exploration. At the time of writing, we have identified the following
  questions that could inspire new innovations:

  \paragraph{Are the developed ideas independent of the underlying
      representation?} This thesis was conducted during a period of rapid advancements in both NeRFs
    and 3DGS. Consequently, methods like \conerf were designed for NeRFs and may
    require adaptation to work with other representations. For instance, NeRFs
    rely on Fourier-transformed input coordinates to generate high-frequency
    details, enabling coarse manipulation of the field by modulating corresponding
    frequencies. In contrast, spatial encodings like hash grids from
    InstantNGP~\cite{mueller2022instant} exhibit high expressivity but are
    challenging to control due to the spatial decoherence of field attributes.
    This discrepancy presents an intriguing avenue for research, especially with
    progress from works like Neuralangelo~\cite{li2023neuralangelo}. Moreover,
    \textcite{yu2024cogs} demonstrated that with proper modifications, \conerf can
    induce controllability in 3DGS as well.

  \paragraph{What other control modalities can be explored?}
    We have explored various control modalities, including sparse
    annotation-induced control in \conerf, spatial texture control with
    \blendfields, lighting control with \lumigauss, and computational adaptability
    with \clog. However, this list is far from exhaustive. Emerging research
    directions include generating assets with text
    prompts~\cite{tang2024dreamgaussian}, human control via parametric
    models~\cite{kocabas2024hugs}, shadow prediction~\cite{solnerf}, and deciding
    between NeRFs and Gaussians for specific 3D volumes~\cite{he2024nerfs}. This
    thesis represents a starting point in a vast and rapidly evolving domain of
    radiance field control.

  \paragraph{Can the introduced concepts extend to generative modeling?}
    Generative modeling for NeRFs and 3DGS currently demands substantial
    computational resources to train a single model~\cite{tang2024dreamgaussian}.
    Most existing approaches borrow from image-based techniques, using text
    prompts for control. A key question is whether our contributions can be
    integrated into such workflows. While no comprehensive solution exists yet, we
    hypothesize that our methods could be applied to radiance fields using a score
    distillation objective~\cite{poole2023dreamfusion}, where both the radiance
    field and the diffusion process are conditioned on control signals. Realizing
    this exciting prospect could unlock novel applications in art and gaming
    industries.
